\section{Contraction of the normalized activity}
\label{subsec:rg-activity-contraction}

On a small polymer the ordinary quadratic kernel now vanishes
at zero connection. The vacuum is also even in the connection,
so its first derivative vanishes there. The analytic
disk gives an additional gain for precisely these two
terms; every other term is controlled directly by its
Grassmann degree.

\begin{proposition}\label{prop:rg-activity-contraction}
Assume \eqref{eq:rg-adjustment-charge-condition},
\eqref{eq:rg-contraction-parameters}, and \eqref{eq:rg-contraction-polymer-parameters}.
For $\mathscr E\in K_k^{x,\C}(\ell)$ set
$E=(R\mathscr E)^{\mathrm{nat}}$.
The linear map $\mathscr E\mapsto\mathcal L E$ takes
$K_k^{x,\C}(\ell)$ into $K_{k+1}^{x,\C}(\ell/L)$,
preserves the real subspaces, and satisfies
\begin{equation}\label{eq:rg-contraction-norm}
 \|\mathcal L E\|_{k+1}^x
       \le c_{\mathrm{fl}}L^{-1/4+3\varepsilon/2}\|\mathscr E\|_k^x
       \le c_{\mathrm{fl}}L^{-1/4+2\varepsilon}\|\mathscr E\|_k^x,
       \qquad c_{\mathrm{fl}}=61c_{\mathrm{rb}}.
\end{equation}
For the weights $w_j^{\mathrm E}=e_j^{1/4-7\varepsilon}$,
\begin{equation}\label{eq:rg-contraction-window}
 \frac{\|\mathcal L E\|_{k+1}^x}{w_{k+1}^{\mathrm E}}
       \le c_{\mathrm{fl}}L^{-3/8+11\varepsilon/2}
                       \frac{\|\mathscr E\|_k^x}{w_k^{\mathrm E}}.
\end{equation}
Restriction to the zero external argument commutes with
the map.
\end{proposition}

\begin{proof}
Write
\[
 N(X)=\sup_{\widetilde R_k\times\mathfrak X_k}
                                      \|E(X;\cdot,\cdot)\|_{h_k}.
\]
Extraction and adjustment give
\begin{equation}\label{eq:rg-adjusted-polymer-bound}
 N(X)\le
 \begin{cases}
 20\|\mathscr E\|_k^x\exp(-\kappa d_M(X)),&X\text{ small},\\
 \|\mathscr E\|_k^x\exp(-\kappa d_M(X)),&X\text{ large}.
 \end{cases}
\end{equation}
For a small polymer this uses
$\|R\mathscr E(X;\mathscr A,\mathscr C)\|_{h_k}
\le\|\mathscr E(X;\mathscr A,\mathscr C)\|_{h_k}
 +\|\mathscr E(X;0,0)\|_{h_k}$ and
\eqref{eq:rg-adjustment-norm}. For a large polymer both
operations leave the activity unchanged.

Fix a small $X$ and a coarse connection pair, and use
the global representatives above. Put $\rho=L^{7/4}$.
The inverse image of the open set $R_k\times\mathfrak X_k$
under $t\mapsto(t\mathscr A^{\mathrm b},t\mathscr C^{\mathrm b})$
contains the closed disk $|t|\le\rho$, hence a slightly
larger disk. The functions
\[
 f_{nm}(t)=E_{nm}(X;t\mathscr A^{\mathrm b},
                              t\mathscr C^{\mathrm b})
\]
are holomorphic there, with values in a finite-dimensional
kernel space. Their weighted norms obey
\[
 \frac{h_{k,1}^nh_{k,2}^m}{n!m!}\|f_{nm}(t)\|_1
                                      \le N(X)\qquad(|t|\le\rho).
\]
Here $\|\cdot\|_1$ is the sum of kernel moduli with their
index weights.

For the vacuum, \eqref{eq:rg-adjusted-kernel-normalization}
gives $f_{00}(0)=0$. Charge conjugation fixes a scalar
kernel and reverses both connections, so
$f_{00}(-t)=f_{00}(t)$ and $f_{00}'(0)=0$.
Cauchy's formula gives
\[
 f_{00}(1)=\frac1{2\pi i}\oint_{|t|=\rho}
                                   \frac{f_{00}(t)}{t^2(t-1)}\,\dd t.
\]
The vacuum equality in
Proposition~\ref{prop:rg-small-polymer-kernel-comparison}
and $\rho\ge2$ therefore give
\begin{equation}\label{eq:rg-vacuum-cauchy-gain}
 \|(S_LE)^{[0,0]}(X;\mathscr A_L,\mathscr C_L)\|_{h_{k+1}}
              \le\frac{N(X)}{\rho(\rho-1)}
              \le2L^{-7/2}N(X).
\end{equation}

For the ordinary quadratic kernel, normalization gives
$f_{20}(0)=0$ as a full antisymmetric kernel. Its balanced
matrix components vanish by
\eqref{eq:rg-adjusted-kernel-normalization}; their
oppositely ordered components vanish by antisymmetry,
and the two equally barred or equally unbarred components
vanish by balance. Hence
\[
 f_{20}(1)=\frac1{2\pi i}\oint_{|t|=\rho}
                                    \frac{f_{20}(t)}{t(t-1)}\,\dd t,
 \qquad
 \|f_{20}(1)\|_1\le\frac{\sup_{|t|=\rho}\|f_{20}(t)\|_1}{\rho-1}.
\]
The comparison of kernels costs $\exp(4)$ in this
bidegree. Its scaling exponent is $p(2,0)=9/4$.
Multiplying the last inequality by $h_{k,1}^2/2$ and
using \eqref{eq:rg-degree-scaling} gives
\begin{equation}\label{eq:rg-quadratic-cauchy-gain}
 \|(S_LE)^{[2,0]}(X;\mathscr A_L,\mathscr C_L)\|_{h_{k+1}}
                       \le2\exp(4)L^{-4}N(X).
\end{equation}

For all remaining bidegrees use the original scaled
connection pair, which belongs to the fine domain.
Evenness leaves only $(1,1)$, $(0,2)$, and total degrees
at least $4$. The exponents satisfy
\[
 p(1,1)>\frac{13}{4}-\frac{3\varepsilon}{2},\qquad
 p(0,2)>\frac{13}{4}-\frac{3\varepsilon}{2},\qquad
 p(n,m)\ge n+m.
\]
Consequently
\[
 \begin{split}
 \sum_{(n,m)\ne(0,0),(2,0)}
       \|(S_LE)^{[n,m]}(X;\mathscr A_L,\mathscr C_L)\|_{h_{k+1}}
 &\le L^{-13/4+3\varepsilon/2}
       \sum_{n,m}\|E^{[n,m]}(X;\mathscr A_L,\mathscr C_L)\|_{h_k}\\
 &\le L^{-13/4+3\varepsilon/2}N(X).
 \end{split}
\]
The sum is taken at one connection pair before bounding
it by $N(X)$. The large-$L$ condition ensures
$2L^{-7/2}\le L^{-13/4}$ and
$2\exp(4)L^{-4}\le L^{-13/4}$. Adding the three
degree contributions yields
\begin{equation}\label{eq:rg-small-polymer-contraction}
 \|S_LE(X;\mathscr A_L,\mathscr C_L)\|_{h_{k+1}}
                       \le3L^{-13/4+3\varepsilon/2}N(X).
\end{equation}

We now sum over the fine polymers in a coarse polymer
$Y$. The hypotheses \eqref{eq:rg-contraction-polymer-parameters}
imply
$\kappa-1\ge\kappa_0$ and
$L(\kappa-\kappa_0-2)\ge\kappa$.
The reblocking bound
\eqref{eq:polymer-reblock-sum} at rate $\kappa$
and \eqref{eq:rg-adjusted-polymer-bound} therefore give
\[
 \begin{split}
 \sum_{\substack{X\text{ small}\\\bar X^L=LY}}
       \|S_LE(X;\mathscr A_L,\mathscr C_L)\|_{h_{k+1}}
 &\le60c_{\mathrm{rb}}L^{-1/4+3\varepsilon/2}
       \|\mathscr E\|_k^x\exp(-\kappa d_M(Y)).
 \end{split}
\]
For large $X$, scaling decreases the norm by
\eqref{eq:rg-scaling-norm-identity} and the ratios of
the two weights. Since $d_M(X)>L$,
\[
 \exp(-\kappa d_M(X))
       \le\exp(-L)\exp(-(\kappa-1)d_M(X)).
\]
Use \eqref{eq:polymer-reblock-sum} at rate $\kappa-1$
to obtain
\[
 \sum_{\substack{X\text{ large}\\\bar X^L=LY}}
       \|S_LE(X;\mathscr A_L,\mathscr C_L)\|_{h_{k+1}}
 \le c_{\mathrm{rb}}L^3\exp(-L)
             \|\mathscr E\|_k^x\exp(-\kappa d_M(Y)).
\]
As $L^3\exp(-L)\le L^{-1/4}$, these estimates prove
\eqref{eq:rg-contraction-norm} after taking the two
suprema in the activity norm. Division by
$w_{k+1}^{\mathrm E}=L^{1/8-7\varepsilon/2}w_k^{\mathrm E}$
gives the relative factor $c_{\mathrm{fl}}L^{-3/8+5\varepsilon}$,
which is at most the factor in \eqref{eq:rg-contraction-window}.
The additional parameter choice
$c_{\mathrm{fl}}L^{-3/8+6\bar\varepsilon}\le1/16$
makes its relative factor at most $1/16$.

The map is linear. Its kernels are finite sums of
holomorphic fine kernels composed with the complex
linear connection scaling. The locality was proved
after \eqref{eq:rg-reblocking-definition}.
For lattice covariance a coarse group element
$g'=(S,a')$ has the fine representative $g=(S,La')$.
It gives a bijection of the polymers in the two
reblocking sums, satisfies
$(g'\mathscr A)_L=g\mathscr A_L$, and intertwines the
spin actions by $a_{g'}S_L=S_La_g$.
The same identity holds for $\mathscr C$.
Thus covariance of each fine activity gives covariance
of its sum. For the real gauges, every fine bond within
distance $1$ of $X$ scales to a bond within distance
$1/L\le1$ of $Y$, while
$G_\lambda S_L=S_LG_{\lambda_L}$ follows from
$e_k\lambda_L(c)=e_{k+1}\lambda(c/L)$.
Independent real gauges of the two arguments therefore
pass through each term. Finally, scaling has real
coefficients, preserves the bar index, and commutes with
charge conjugation and the antilinear reality map.
These facts prove membership in both stated activity
spaces. Since extraction and adjustment commute with
restriction and $0_L=0$, reblocking also commutes with
restriction to the zero external argument.
\end{proof}

The reblocked sum retains the evaluated normalized
activity. Indeed, each fine polymer belongs to exactly
one reblocking class, so
\begin{equation}\label{eq:rg-reblocked-sum}
 \sum_Y(\mathcal LE)(Y;\mathscr A,\mathscr C)
      =S_L\sum_X E(X;\mathscr A_L,\mathscr C_L).
\end{equation}
If the coarse generators are evaluated at fields
with $\chi'=\delta_{\alpha,\mathscr A+\mathscr C}\psi'$
and its barred counterpart, their scaled fields satisfy
the fine relations. This follows from
$d_X(x/L,x'/L)=d_X(x,x')/L$ and
$U_{(\mathscr A+\mathscr C)_L}(x,x')
 =U_{\mathscr A+\mathscr C}(x/L,x'/L)$:
the path bijection and
$e_k\eta_kL^{-1/2}=e_{k+1}\eta_{k+1}$ give the latter.
Substitution in the difference quotient produces
$L^{-1-\alpha}\chi'(x/L,x'/L)$, and the barred quotient
uses the inverse transport. Adjustment consequently
has its original value under this evaluation. Thus the
evaluation of \eqref{eq:rg-reblocked-sum} is also the
evaluation of
\[
 S_L\sum_X(R\mathscr E)(X;\mathscr A_L,\mathscr C_L).
\]
\begin{proof}[Proof of Theorem~\ref{thm:main}]
Proposition~\ref{prop:rg-scalar-counterterms} gives the scalar
counterterms and their bounds. Equations
\eqref{eq:rg-exact-extraction-identity} and
\eqref{eq:rg-remainder-properties} give the exact extraction and
the integrated normalization. The shortest-chain construction
\eqref{eq:rg-chain-homomorphism} preserves evaluated values;
\eqref{eq:rg-adjusted-kernel-normalization} and
\eqref{eq:rg-adjustment-norm} give the full kernel normalization
and its bound. Proposition~\ref{prop:rg-activity-contraction}
proves the contraction and the relative estimate. The covariance
and restriction statements were proved with each map.
\end{proof}
